\documentclass[english, parskip=full]{scrartcl}
\usepackage[utf8]{inputenc}  % Kodierung der Datei
\usepackage[english]{babel}  % Sprache auf Deutsch 
\usepackage{color}
\usepackage{graphicx, gensymb}
\usepackage{amsfonts, cancel, amssymb}
\usepackage{amsmath, amsthm, array, esint}
\usepackage{booktabs, multirow}  % schönere Tabellen
\usepackage{caption}  % Anpassung der Captions
\usepackage{scrlayer-scrpage}    % Manipulation des Seitenstils
\pagestyle{scrheadings}  % Stil für die Seite setzen
\automark{section}  % Abschnittsnamen als Seitenbeschriftung 
\setheadsepline{.5pt}    % Kopzeile durch Linie abtrennen
\usepackage[hidelinks]{hyperref}  % Links und weitere PDF-Features
\usepackage{typearea, tikz, pgffor, verbatim}
\usetikzlibrary{calc, arrows.meta,arrows, graphs, quotes, positioning,angles}
\usepackage[cache=false, outputdir=build]{minted}
\usemintedstyle{vs}
\usepackage{placeins} % provides \FloatBarrier

\definecolor{bg}{rgb}{0.9,0.9,0.9}
\newcommand\numberthis{\addtocounter{equation}{1}\tag{\theequation}}
\usepackage{svg}
\usepackage{siunitx}
\usepackage{physics}
\usepackage{tensor}
\usepackage[compat=1.1.0]{tikz-feynman}
\renewcommand{\bra}{}
\newcommand{\kom}{}
\newcommand{\brak}{}
\renewcommand{\braket}{}
\newcommand{\intx}{}
\newcommand{\intr}{}
\newcommand{\kla}{}
\newcommand{\klam}{}
\newcommand{\klammer}{}
\newcommand{\oper}{}
\newcommand{\tecxt}{}
\newcommand{\Bra}{}
\newcommand{\Ket}{}
\renewcommand{\i}{\text{i}}
\newcommand{\e}{\text{e}}
\newcommand*{\Fourier}{\textsc{Fourier}\ }
\newcommand*{\F}{\mathcal{F}}
\newcommand*{\sinc}{\text{sinc\,}}
\newcommand{\conv}{\mathop{\scalebox{3.5}{\raisebox{-0.38ex}{\makebox[0.5\width][c]{$\ast$}}}}\limits}

\newcommand*{\titel}{Generalized Numerical Analysis of Fundamental Frequencies}
\newcommand*{\autor}{Joachim Schwardt}

\hypersetup{pdfauthor={\autor}, pdftitle={\titel}}

%\titlehead{titlehead}
\subject{Frequency Analysis}
\title{\titel}
\author{\autor}
\date{\today}

\begin{document}

%\begin{titlepage}
\maketitle  % Titel setzen
\tableofcontents  % Inhaltsverzeichnis setzen
%\end{titlepage}

\section{Overview}
In this section we want to introduce the problem and illustrate how we will solve it without going too deep into the technical details.
Consider a time-series, i.e. a sequence of $N$ complex numbers
\begin{align}
Z_N &:= \klammer\left\{ z_n \, \vert \, n=0,\dots,N-1 \right\}
\end{align}
which we will also refer to as a \emph{signal}.
One could in principle also consider higher-dimensional variants.
But if the components have a sufficient coupling, one can just extract all frequencies from either one of them, and otherwise the problem separates into one-dimensional problems anyway.
As such there might be additional complications, but the general approach should still remain the same.
\\
We want to find the \emph{fundamental frequencies} of $Z_N$.
For this we define a model of our signal as
\begin{align}
z_n^\tecxt\text{M} &:= \sum_{k=0}^K c_k \e^{2\pi\i \nu_k n}.
\end{align}
Here the $c_k$ are the amplitudes and phases of the frequency $\nu_k$ in our signal, and $K$ is total number of frequencies.
The superscript M is to remind one that this is not the measured signal, but rather a reconstruction using Fourier-components.
For this approach to work with only few such components, the Fourier-transformation of the $z_n$ will generally have to contain isolated peaks as only a few frequencies.
A slightly more precise, while still somewhat ambiguous statement of this idea is that the Fourier-components
\begin{align}
F_j &:= \sum_{n=0}^{N-1} z_n \e^{-2\pi\i \frac{j}{N} n}
\end{align}
should only have large amplitudes for a number of indices $j\in\{0,\dots,N-1\}$ that is in a certain sense small compared to the total number of components $N$.\\
So far, this setup is rather simple.
The generalization, and ultimately improved frequency analysis, stems from the idea of weighting the signal using weights $w_n$ drawn from a weighting or window function $w(x)$ defined on the interval $[0,1]$.
This window function should drop of to zero at the edges, to ensure that concatenating the weighted signal $w_nz_n$ with itself repeatedly results in as smooth of a periodic signal as possible.
To give an example, consider 
\begin{align}
z_n := \sin\kla\left( 2\pi \frac{n}{N} \cdot \frac{3}{2} \right).
\end{align}



%\begin{thebibliography}{9}
%\bibitem{example} description, \url{https://www.exmapleurl}, day.\,month~2021.
%\end{thebibliography}

\end{document}